\section{Introduction}
Stochastic gradient Langevin dynamics (SGLD) replaces the full posterior gradient by a minibatch estimate to make Bayesian computation feasible on large datasets \citep{welling2011}. The resulting subsampling noise is distinct from the injected Langevin noise. First-order control variates anchor the stochastic gradient at a reference point, reducing its variability near that anchor \citep{baker2017}. A second-order control variate additionally cancels the local linear variation of each observation's gradient. The resulting estimator and its theoretical analysis are developed in Sections~\ref{sec:method} and~\ref{sec:theory}.

Our experimental question is when this additional variance reduction improves the accuracy--cost tradeoff of posterior sampling. We focus on a fixed inverse Hessian of the negative log posterior at the anchor, including prior curvature and declared numerical regularization. This separates the gradient-estimator comparison from online adaptive-metric dynamics. Full-gradient unadjusted Langevin sampling (ULA) supplies a diagnostic control for the contribution of stochastic gradients, without being treated as a zero-error posterior sampler.

The empirical evidence has three complementary parts. First, frozen German Credit and logistic-regression comparisons establish accuracy improvements as well as unmet efficiency criteria. Second, an original Bank Marketing failure motivates implementation checks and controlled observation-level interventions. These identify a high-leverage Taylor-remainder mechanism; a transformed-target study then tests the optimized accuracy--cost frontier with independent validation. Third, a nested synthetic study examines dataset-size scaling of both gradient noise and tuned wall-clock cost. A focused Gaussian experiment verifies algebraic exactness and the geometry correction separately from practical speedup.

Together these experiments support conditional computational benefits. They also explain why noise reduction can fail to deliver an advantage: rare large updates, discretization error, grid resolution, competing parameter choices and shared setup costs all matter. We retain unfavorable results and distinguish original primary comparisons from historically motivated extensions and post-result sensitivity analyses.

\section{Related Works}\label{sec:related}
SGLD was introduced by \citet{welling2011}. The control-variate construction of \citet{baker2017} is the first-order baseline extended here. Inaccurate-gradient Langevin bounds provide context for the theoretical analysis \citep{dalalyan2017}, and Hessian--vector products provide a general implementation route for curvature corrections \citep{pearlmutter1994}. The experiments use cached likelihood curvature where practical, while measuring the full posterior Hessian required for the sampling metric.

Adaptive preconditioning is a distinct approach. \citet{li2016psgld} combine an evolving gradient-based diagonal metric with SGLD. Our primary comparisons instead share a fixed full-Hessian metric across CV1, CV2 and full-gradient ULA, holding posterior geometry constant to isolate stochastic-gradient variance reduction. CV2's cached anchor likelihood Hessian can be reused for this metric after adding prior curvature and the declared regularization; all methods are charged for their required metric preparation. Combining control variates with an adaptive metric introduces additional choices about accumulator inputs, initialization and adaptive drift corrections. Those choices are outside the primary comparison. Earlier adaptive controls are retained for completeness, without a general conclusion about Li et al.'s method. Full-data Hamiltonian Monte Carlo (HMC) supplies approximate reference draws \citep{neal2012hmc}; their quality is assessed using rank-based convergence and effective-sample-size diagnostics \citep{vehtari2021rank}.
